


\documentclass{scrartcl}\usepackage[]{graphicx}\usepackage[]{color}
%% maxwidth is the original width if it is less than linewidth
%% otherwise use linewidth (to make sure the graphics do not exceed the margin)
\makeatletter
\def\maxwidth{ %
  \ifdim\Gin@nat@width>\linewidth
    \linewidth
  \else
    \Gin@nat@width
  \fi
}
\makeatother

\definecolor{fgcolor}{rgb}{0.345, 0.345, 0.345}
\newcommand{\hlnum}[1]{\textcolor[rgb]{0.686,0.059,0.569}{#1}}%
\newcommand{\hlstr}[1]{\textcolor[rgb]{0.192,0.494,0.8}{#1}}%
\newcommand{\hlcom}[1]{\textcolor[rgb]{0.678,0.584,0.686}{\textit{#1}}}%
\newcommand{\hlopt}[1]{\textcolor[rgb]{0,0,0}{#1}}%
\newcommand{\hlstd}[1]{\textcolor[rgb]{0.345,0.345,0.345}{#1}}%
\newcommand{\hlkwa}[1]{\textcolor[rgb]{0.161,0.373,0.58}{\textbf{#1}}}%
\newcommand{\hlkwb}[1]{\textcolor[rgb]{0.69,0.353,0.396}{#1}}%
\newcommand{\hlkwc}[1]{\textcolor[rgb]{0.333,0.667,0.333}{#1}}%
\newcommand{\hlkwd}[1]{\textcolor[rgb]{0.737,0.353,0.396}{\textbf{#1}}}%
\let\hlipl\hlkwb

\usepackage{framed}
\makeatletter
\newenvironment{kframe}{%
 \def\at@end@of@kframe{}%
 \ifinner\ifhmode%
  \def\at@end@of@kframe{\end{minipage}}%
  \begin{minipage}{\columnwidth}%
 \fi\fi%
 \def\FrameCommand##1{\hskip\@totalleftmargin \hskip-\fboxsep
 \colorbox{shadecolor}{##1}\hskip-\fboxsep
     % There is no \\@totalrightmargin, so:
     \hskip-\linewidth \hskip-\@totalleftmargin \hskip\columnwidth}%
 \MakeFramed {\advance\hsize-\width
   \@totalleftmargin\z@ \linewidth\hsize
   \@setminipage}}%
 {\par\unskip\endMakeFramed%
 \at@end@of@kframe}
\makeatother

\definecolor{shadecolor}{rgb}{.97, .97, .97}
\definecolor{messagecolor}{rgb}{0, 0, 0}
\definecolor{warningcolor}{rgb}{1, 0, 1}
\definecolor{errorcolor}{rgb}{1, 0, 0}
\newenvironment{knitrout}{}{} % an empty environment to be redefined in TeX

\usepackage{alltt}
\usepackage[margin=.5in,top=.5in,bottom=.3in,includehead,includefoot]{geometry}
\renewcommand{\arraystretch}{1.5}

%\parindent=0pt
\parskip=3pt


%\usepackage{mathpple}
%\usepackage{fullpage}
\usepackage{alltt}
\usepackage{sfsect}
\usepackage{longtable}
\usepackage[yyyymmdd,hhmmss]{datetime}


\def\WW#1{\href{https://mathwebwork.cs.calvin.edu/webwork2/Math252-S19-Pruim/WW#1}{WW#1}}

\newif\ifanswers
\answerstrue
\answersfalse

\def\blank#1{%
\ifanswers
\underline{{\large \sc #1}}
\else
\underline{\phantom{\Large \ \ \ \sc #1 \ \ \ }}
\fi
}

\def\dd#1{\medskip\noindent\textbf{\large #1}}
\def\be{\begin{enumerate}}
\def\ee{\end{enumerate}}
\def\bi{\begin{itemize}}
\def\ei{\end{itemize}}

\def\hwn#1{\relax}
\def\hwn#1{\ {\tiny (#1)}}
\def\hwn#1{$_{\mbox{\ \tiny #1}}$}
\def\hwn#1{~{\tiny #1}}

\usepackage[normalem]{ulem}
\usepackage{ifthen}
\usepackage{fancyhdr}
\usepackage[colorlinks = true, urlcolor = blue]{hyperref}
\include{myabbrev}


\renewcommand{\labelenumi}{\textbf{\arabic{enumi}}}
\newcommand{\myunderline}[1]{\uline{\phantom{Ay}#1\phantom{Ay}}}

\newcommand{\helper}[2][\relax]{(#2)\hwn{#1}}
\newcommand{\extra}[2][\relax]{[#2]\hwn{#1}}
\newcommand{\opt}[2][\relax]{\{#2\}\hwn{#1}}
\newcommand{\mand}[2][\relax]{{\bf \large #2}\hwn{#1}}
\IfFileExists{upquote.sty}{\usepackage{upquote}}{}
\begin{document}


\pagestyle{fancy}

\chead{}
\lhead{Math 252 -- Spring 2019}
\rhead{\thepage/\pageref{end}}
\lfoot{}
\cfoot{Created \today\ --- See web site for most current version.}
\rfoot{}

\begin{center}
  \large
%  Mathematics \& Statistics 156
%
%  Fall 2016
%
\subsection*{Problem Sets After Test 2}

\end{center}


\textit{Only turn in problems that are \textbf{not} bracketed.}
Bracketed problems are additional problems you can look at.
Round brackets indicate problems that may help you with problems that
are assigned; square brackets are additional problems on material
that you should know, but you are not required to write up solutions;
curly brackets are truly optional and may contain extra nuggets that you
will not be required to know but may be interested in.

Additional assignments will be filled in over time.

\bigskip

\begin{center}

\iftrue

\begin{longtable}{|c|l|}
  \hline
  notation & meaning \\
  \hline
  unbracketed & assigned problem -- turn these in for grading \\
  $()$ & helper/warm-up problem  \\
  $[]$ & additional problems (you are responsible for content, but don't turn them in) \\
  \{\} & covers optional material \\
  \hline
\end{longtable}
\fi

\bigskip


\begin{longtable}{|c|c|c|p{4.3in}|}
  \hline
  Due & Task & Source & Problems \\
  \hline\hline
  \endhead
  Fri 4/12 & PS 9 & Rosen 13.1 &
   \mand[sentences]{1ad}
	 \extra{2}
	 \helper{3}
	 \mand{4}
	 \extra{5}
	 \mand{6a}
	 \mand{9}
	 \helper{13}
	 \mand{14bd}
	 \extra{15}

	 % \mand[types of grammars]{19}
	 % \mand[palindromes]{20}
	 % \mand[union]{21a}
	 % \mand[derivation trees]{24ab}

	 % \mand[BNF]{28}
	 % \mand[BNF]{31ab}
	 % \mand[EBNF]{34}
   \\
	   & & Rosen 13.3 &
	   \helper[testing strings]{11}
	   \mand[testing strings]{12}
	   \mand[what langauge?]{16--17}
	   \mand[create DFA]{25--26}
   \\
   \hline
	   Mon 4/15 & PS 10 & Rosen 13.3 &
	   \mand[what langauge?]{19--20}
	   \mand[create DFA]{32}
	   \mand[create DFA]{34}
	   \mand[what langauge?]{45--46}
	   \mand[NFA to DFA]{52--54}
   \\
	   \hline
	   Wed 4/24 & PS 11 & Rosen 13.4 &
	   \extra{1}
	   \mand{2}
	   \mand[regex]{6}
	   \extra[regex]{7}
	   \mand[Kleene's Thm]{12--13}
	   \mand[lonely start state]{20}
   \\
	   \hline
	   Mon 4/29 & PS 12 & Rosen 13.4 &
	   \mand[grammar $\to$ NFA]{14}
	   \mand[NFA $\to$ grammar]{15--16}
	   \mand[pumping lemma]{22, 24, 25}
   \\
	   \hline
	   	  Wed 4/31 & PS 13 & &
	  Prepare for Test 3.  Here are some problems you might look at.
	  \\
   & & Rosen 13.RQ &
   \extra{1--2}
   \extra{4}
   \extra{5a}
   \extra{6a}
   \extra{7--12}
   \extra{15}
   \\
   & & Rosen 13.SE &
   \extra{3}
   \extra{7}
   \extra{8}
   \extra{9}
   \extra{19--21}
   \extra{23--25}
   \extra{26--27}
   	   \\
	   \hline
	  ** & PS 14 & Rosen 13.5 &
	  \extra[trace TM]{1--5}
	  \extra[design TM]{6--12}
	   \\
	   \hline
   \end{longtable}
   \end{center}
   \end{document}



	   \\
	   \hline

\end{longtable}
\end{center}
\label{end}
\end{document}
\end{document}

